\documentclass[
  jou,
  floatsintext,
  longtable,
  a4paper,
  nolmodern,
  notxfonts,
  notimes,
  12pt,
  colorlinks=true,linkcolor=blue,citecolor=blue,urlcolor=blue]{apa7}

\usepackage{amsmath}
\usepackage{amssymb}



\usepackage[bidi=default]{babel}
\babelprovide[main,import]{spanish}


% get rid of language-specific shorthands (see #6817):
\let\LanguageShortHands\languageshorthands
\def\languageshorthands#1{}

\RequirePackage{longtable}
\RequirePackage{threeparttablex}

\makeatletter
\renewcommand{\paragraph}{\@startsection{paragraph}{4}{\parindent}%
	{0\baselineskip \@plus 0.2ex \@minus 0.2ex}%
	{-.5em}%
	{\normalfont\normalsize\bfseries\typesectitle}}

\renewcommand{\subparagraph}[1]{\@startsection{subparagraph}{5}{0.5em}%
	{0\baselineskip \@plus 0.2ex \@minus 0.2ex}%
	{-\z@\relax}%
	{\normalfont\normalsize\bfseries\itshape\hspace{\parindent}{#1}\textit{\addperi}}{\relax}}
\makeatother




\usepackage{longtable, booktabs, multirow, multicol, colortbl, hhline, caption, array, float, xpatch}
\usepackage{subcaption}


\renewcommand\thesubfigure{\Alph{subfigure}}
\setcounter{topnumber}{2}
\setcounter{bottomnumber}{2}
\setcounter{totalnumber}{4}
\renewcommand{\topfraction}{0.85}
\renewcommand{\bottomfraction}{0.85}
\renewcommand{\textfraction}{0.15}
\renewcommand{\floatpagefraction}{0.7}

\usepackage{tcolorbox}
\tcbuselibrary{listings,theorems, breakable, skins}
\usepackage{fontawesome5}

\definecolor{quarto-callout-color}{HTML}{909090}
\definecolor{quarto-callout-note-color}{HTML}{0758E5}
\definecolor{quarto-callout-important-color}{HTML}{CC1914}
\definecolor{quarto-callout-warning-color}{HTML}{EB9113}
\definecolor{quarto-callout-tip-color}{HTML}{00A047}
\definecolor{quarto-callout-caution-color}{HTML}{FC5300}
\definecolor{quarto-callout-color-frame}{HTML}{ACACAC}
\definecolor{quarto-callout-note-color-frame}{HTML}{4582EC}
\definecolor{quarto-callout-important-color-frame}{HTML}{D9534F}
\definecolor{quarto-callout-warning-color-frame}{HTML}{F0AD4E}
\definecolor{quarto-callout-tip-color-frame}{HTML}{02B875}
\definecolor{quarto-callout-caution-color-frame}{HTML}{FD7E14}

%\newlength\Oldarrayrulewidth
%\newlength\Oldtabcolsep


\usepackage{hyperref}



\usepackage{color}
\usepackage{fancyvrb}
\newcommand{\VerbBar}{|}
\newcommand{\VERB}{\Verb[commandchars=\\\{\}]}
\DefineVerbatimEnvironment{Highlighting}{Verbatim}{commandchars=\\\{\}}
% Add ',fontsize=\small' for more characters per line
\usepackage{framed}
\definecolor{shadecolor}{RGB}{241,243,245}
\newenvironment{Shaded}{\begin{snugshade}}{\end{snugshade}}
\newcommand{\AlertTok}[1]{\textcolor[rgb]{0.68,0.00,0.00}{#1}}
\newcommand{\AnnotationTok}[1]{\textcolor[rgb]{0.37,0.37,0.37}{#1}}
\newcommand{\AttributeTok}[1]{\textcolor[rgb]{0.40,0.45,0.13}{#1}}
\newcommand{\BaseNTok}[1]{\textcolor[rgb]{0.68,0.00,0.00}{#1}}
\newcommand{\BuiltInTok}[1]{\textcolor[rgb]{0.00,0.23,0.31}{#1}}
\newcommand{\CharTok}[1]{\textcolor[rgb]{0.13,0.47,0.30}{#1}}
\newcommand{\CommentTok}[1]{\textcolor[rgb]{0.37,0.37,0.37}{#1}}
\newcommand{\CommentVarTok}[1]{\textcolor[rgb]{0.37,0.37,0.37}{\textit{#1}}}
\newcommand{\ConstantTok}[1]{\textcolor[rgb]{0.56,0.35,0.01}{#1}}
\newcommand{\ControlFlowTok}[1]{\textcolor[rgb]{0.00,0.23,0.31}{\textbf{#1}}}
\newcommand{\DataTypeTok}[1]{\textcolor[rgb]{0.68,0.00,0.00}{#1}}
\newcommand{\DecValTok}[1]{\textcolor[rgb]{0.68,0.00,0.00}{#1}}
\newcommand{\DocumentationTok}[1]{\textcolor[rgb]{0.37,0.37,0.37}{\textit{#1}}}
\newcommand{\ErrorTok}[1]{\textcolor[rgb]{0.68,0.00,0.00}{#1}}
\newcommand{\ExtensionTok}[1]{\textcolor[rgb]{0.00,0.23,0.31}{#1}}
\newcommand{\FloatTok}[1]{\textcolor[rgb]{0.68,0.00,0.00}{#1}}
\newcommand{\FunctionTok}[1]{\textcolor[rgb]{0.28,0.35,0.67}{#1}}
\newcommand{\ImportTok}[1]{\textcolor[rgb]{0.00,0.46,0.62}{#1}}
\newcommand{\InformationTok}[1]{\textcolor[rgb]{0.37,0.37,0.37}{#1}}
\newcommand{\KeywordTok}[1]{\textcolor[rgb]{0.00,0.23,0.31}{\textbf{#1}}}
\newcommand{\NormalTok}[1]{\textcolor[rgb]{0.00,0.23,0.31}{#1}}
\newcommand{\OperatorTok}[1]{\textcolor[rgb]{0.37,0.37,0.37}{#1}}
\newcommand{\OtherTok}[1]{\textcolor[rgb]{0.00,0.23,0.31}{#1}}
\newcommand{\PreprocessorTok}[1]{\textcolor[rgb]{0.68,0.00,0.00}{#1}}
\newcommand{\RegionMarkerTok}[1]{\textcolor[rgb]{0.00,0.23,0.31}{#1}}
\newcommand{\SpecialCharTok}[1]{\textcolor[rgb]{0.37,0.37,0.37}{#1}}
\newcommand{\SpecialStringTok}[1]{\textcolor[rgb]{0.13,0.47,0.30}{#1}}
\newcommand{\StringTok}[1]{\textcolor[rgb]{0.13,0.47,0.30}{#1}}
\newcommand{\VariableTok}[1]{\textcolor[rgb]{0.07,0.07,0.07}{#1}}
\newcommand{\VerbatimStringTok}[1]{\textcolor[rgb]{0.13,0.47,0.30}{#1}}
\newcommand{\WarningTok}[1]{\textcolor[rgb]{0.37,0.37,0.37}{\textit{#1}}}

\providecommand{\tightlist}{%
  \setlength{\itemsep}{0pt}\setlength{\parskip}{0pt}}
\usepackage{longtable,booktabs,array}
\newcounter{none} % for unnumbered tables
\usepackage{calc} % for calculating minipage widths
% Correct order of tables after \paragraph or \subparagraph
\usepackage{etoolbox}
\makeatletter
\patchcmd\longtable{\par}{\if@noskipsec\mbox{}\fi\par}{}{}
\makeatother
% Allow footnotes in longtable head/foot
\IfFileExists{footnotehyper.sty}{\usepackage{footnotehyper}}{\usepackage{footnote}}
\makesavenoteenv{longtable}

\usepackage{graphicx}
\makeatletter
\newsavebox\pandoc@box
\newcommand*\pandocbounded[1]{% scales image to fit in text height/width
  \sbox\pandoc@box{#1}%
  \Gscale@div\@tempa{\textheight}{\dimexpr\ht\pandoc@box+\dp\pandoc@box\relax}%
  \Gscale@div\@tempb{\linewidth}{\wd\pandoc@box}%
  \ifdim\@tempb\p@<\@tempa\p@\let\@tempa\@tempb\fi% select the smaller of both
  \ifdim\@tempa\p@<\p@\scalebox{\@tempa}{\usebox\pandoc@box}%
  \else\usebox{\pandoc@box}%
  \fi%
}
% Set default figure placement to htbp
\def\fps@figure{htbp}
\makeatother







\usepackage{newtx}

\defaultfontfeatures{Scale=MatchLowercase}
\defaultfontfeatures[\rmfamily]{Ligatures=TeX,Scale=1}





\title{Guía completa de instalación de R y RStudio Desktop}


\shorttitle{Guía completa de instalación de R y RStudio Desktop}


\usepackage{etoolbox}



\ccoppy{\textcopyright~2020}



\author{Edison Achalma}



\affiliation{
{Escuela Profesional de Economía, Universidad Nacional de San Cristóbal
de Huamanga}}




\leftheader{Achalma}

\date{2020-06-10}


\abstract{Este abstract será actualizado una vez que se complete el
contenido final del artículo. }

\keywords{keyword1, keyword2}

\authornote{\par{\addORCIDlink{Edison Achalma}{0000-0001-6996-3364}} 
\par{ }
\par{   El autor no tiene conflictos de interés que revelar.    Los
roles de autor se clasificaron utilizando la taxonomía de roles de
colaborador (CRediT; https://credit.niso.org/) de la siguiente
manera:  Edison Achalma:   conceptualización, metodología, análisis
formal, investigación, recursos, curación de
datos, redacción, visualización, supervisión, administración del
proyecto}
\par{La correspondencia relativa a este artículo debe dirigirse a Edison
Achalma, Escuela Profesional de Economía, Universidad Nacional de San
Cristóbal de Huamanga, Portal Independencia N°
57, Ayacucho, AYA 5001, Perú, Email: \href{mailto:elmer.achalma.09@unsch.edu.pe}{elmer.achalma.09@unsch.edu.pe}}
}

\usepackage{pbalance}
% \usepackage{float}
\makeatletter
\let\oldtpt\ThreePartTable
\let\endoldtpt\endThreePartTable
\def\ThreePartTable{\@ifnextchar[\ThreePartTable@i \ThreePartTable@ii}
\def\ThreePartTable@i[#1]{\begin{figure}[!htbp]
\onecolumn
\begin{minipage}{0.485\textwidth}
\oldtpt[#1]
}
\def\ThreePartTable@ii{\begin{figure}[!htbp]
\onecolumn
\begin{minipage}{0.48\textwidth}
\oldtpt
}
\def\endThreePartTable{
\endoldtpt
\end{minipage}
\twocolumn
\end{figure}}
\makeatother


\makeatletter
\let\endoldlt\endlongtable		
\def\endlongtable{
\hline
\endoldlt}
\makeatother

\newenvironment{twocolumntable}% environment name
{% begin code
\begin{table*}[!htbp]%
\onecolumn%
}%
{%
\twocolumn%
\end{table*}%
}% end code

\urlstyle{same}



\hypersetup{
  pdftitle={Gu\'{\i}a completa de instalaci\'{o}n de R y RStudio Desktop},
  pdfauthor={Autor}
}
\makeatletter
\@ifpackageloaded{caption}{}{\usepackage{caption}}
\AtBeginDocument{%
\ifdefined\contentsname
  \renewcommand*\contentsname{Tabla de contenidos}
\else
  \newcommand\contentsname{Tabla de contenidos}
\fi
\ifdefined\listfigurename
  \renewcommand*\listfigurename{Lista de Figuras}
\else
  \newcommand\listfigurename{Lista de Figuras}
\fi
\ifdefined\listtablename
  \renewcommand*\listtablename{Lista de Tablas}
\else
  \newcommand\listtablename{Lista de Tablas}
\fi
\ifdefined\figurename
  \renewcommand*\figurename{Figura}
\else
  \newcommand\figurename{Figura}
\fi
\ifdefined\tablename
  \renewcommand*\tablename{Tabla}
\else
  \newcommand\tablename{Tabla}
\fi
}
\@ifpackageloaded{float}{}{\usepackage{float}}
\floatstyle{ruled}
\@ifundefined{c@chapter}{\newfloat{codelisting}{h}{lop}}{\newfloat{codelisting}{h}{lop}[chapter]}
\floatname{codelisting}{Listado}
\newcommand*\listoflistings{\listof{codelisting}{Lista de Listados}}
\makeatother
\makeatletter
\makeatother
\makeatletter
\@ifpackageloaded{caption}{}{\usepackage{caption}}
\@ifpackageloaded{subcaption}{}{\usepackage{subcaption}}
\makeatother
\makeatletter
\@ifpackageloaded{fontawesome5}{}{\usepackage{fontawesome5}}
\makeatother

% From https://tex.stackexchange.com/a/645996/211326
%%% apa7 doesn't want to add appendix section titles in the toc
%%% let's make it do it
\makeatletter
\xpatchcmd{\appendix}
  {\par}
  {\addcontentsline{toc}{section}{\@currentlabelname}\par}
  {}{}
\makeatother

%% Disable longtable counter
%% https://tex.stackexchange.com/a/248395/211326

\usepackage{etoolbox}

\makeatletter
\patchcmd{\LT@caption}
  {\bgroup}
  {\bgroup\global\LTpatch@captiontrue}
  {}{}
\patchcmd{\longtable}
  {\par}
  {\par\global\LTpatch@captionfalse}
  {}{}
\apptocmd{\endlongtable}
  {\ifLTpatch@caption\else\addtocounter{table}{-1}\fi}
  {}{}
\newif\ifLTpatch@caption
\makeatother

\begin{document}

\maketitle


\hypertarget{toc}{}
\tableofcontents
\newpage
\section[Introduction]{Guía completa de instalación de R y RStudio
Desktop}

\setcounter{secnumdepth}{3}

\setlength\LTleft{0pt}




Referencia técnica independiente sobre la instalación de \textbf{R}
(lenguaje y entorno estadístico) y \textbf{RStudio Desktop} (IDE de
Posit, PBC) en tres sistemas: Kubuntu/Ubuntu (Debian-based), Arch Linux,
y Windows. Basada en la documentación oficial de CRAN
(cran.r-project.org), ArchWiki (wiki.archlinux.org) y Posit Docs
(docs.posit.co / posit.co).

\section{1. Conceptos previos}\label{conceptos-previos}

\subsection{1.1. R vs.~RStudio}\label{r-vs.-rstudio}

\textbf{R} es el lenguaje de programación y el entorno de ejecución para
computación estadística y gráficos, desarrollado por la R Foundation. Se
instala primero porque es el ``motor''.

\textbf{RStudio Desktop} es un IDE (entorno de desarrollo integrado)
desarrollado por \textbf{Posit, PBC} (antes RStudio, PBC) que se ejecuta
\emph{sobre} una instalación de R ya existente. RStudio no incluye R:
siempre hay que instalar R primero y RStudio después, y RStudio detecta
automáticamente la(s) instalación(es) de R presentes en el sistema.

\subsection{1.2. Orden de instalación
recomendado}\label{orden-de-instalaciuxf3n-recomendado}

\begin{enumerate}
\def\labelenumi{\arabic{enumi}.}
\tightlist
\item
  Instalar R (desde el repositorio oficial de tu distribución o desde
  CRAN).
\item
  Verificar que R funciona desde la terminal (\texttt{R\ -\/-version}).
\item
  Instalar RStudio Desktop (paquete binario oficial de Posit).
\item
  Verificar que RStudio detecta la instalación de R.
\end{enumerate}

\subsection{1.3. Compatibilidad de plataformas (RStudio
Desktop)}\label{compatibilidad-de-plataformas-rstudio-desktop}

Es importante saber que \textbf{Posit ya no construye RStudio Desktop
para todas las distribuciones Linux indefinidamente}. Desde la versión
2025.05.0 (``Mariposa Orchid''), Posit dejó de generar binarios de
RStudio Desktop para SLES 15, Ubuntu 20.04 LTS (Focal Fossa) y RHEL 8,
debido a dependencias de software incompatibles. Si usas Kubuntu/Ubuntu,
se recomienda una versión 22.04 LTS o superior para garantizar
compatibilidad con las versiones más recientes de RStudio.

\begin{center}\rule{0.5\linewidth}{0.5pt}\end{center}

\section{2. Instalación de R en Kubuntu /
Ubuntu}\label{instalaciuxf3n-de-r-en-kubuntu-ubuntu}

Kubuntu comparte la base de paquetes de Ubuntu (APT/\texttt{.deb}), por
lo que todo lo siguiente aplica igual a ambas. Existen dos rutas
oficiales: el repositorio de Ubuntu (más simple, versión puede no ser la
más reciente) o el repositorio de \textbf{CRAN} (versión más
actualizada, recomendado para trabajo serio en R).

\subsection{2.1. Opción A --- Repositorio de Ubuntu (rápida, versión
potencialmente
desactualizada)}\label{opciuxf3n-a-repositorio-de-ubuntu-ruxe1pida-versiuxf3n-potencialmente-desactualizada}

\begin{Shaded}
\begin{Highlighting}[]
\FunctionTok{sudo}\NormalTok{ apt update}
\FunctionTok{sudo}\NormalTok{ apt install }\AttributeTok{{-}{-}no{-}install{-}recommends}\NormalTok{ r{-}base r{-}base{-}dev}
\end{Highlighting}
\end{Shaded}

Esta vía es la más simple, pero la versión de R que entrega depende del
ciclo de paquetes de tu versión de Ubuntu/Kubuntu y puede quedar
rezagada respecto a la última versión estable de R.

\subsection{2.2. Opción B --- Repositorio oficial de CRAN
(recomendada)}\label{opciuxf3n-b-repositorio-oficial-de-cran-recomendada}

CRAN mantiene un repositorio APT específico para Ubuntu con builds
actualizados de R 4.x, compilados por Michael Rutter y mantenido junto
con Dirk Eddelbuettel. Estos son los pasos oficiales (README breve de
CRAN):

\textbf{Paso 1 --- Actualizar índices e instalar paquetes auxiliares}

\begin{Shaded}
\begin{Highlighting}[]
\FunctionTok{sudo}\NormalTok{ apt update }\AttributeTok{{-}qq}
\FunctionTok{sudo}\NormalTok{ apt install }\AttributeTok{{-}{-}no{-}install{-}recommends}\NormalTok{ software{-}properties{-}common dirmngr}
\end{Highlighting}
\end{Shaded}

\begin{quote}
En versiones recientes de Ubuntu/Kubuntu,
\texttt{software-properties-common} y \texttt{dirmngr} suelen venir
preinstalados; el comando no hará daño si ya están presentes.
\end{quote}

\textbf{Paso 2 --- Importar la clave de firma del repositorio CRAN}

\begin{Shaded}
\begin{Highlighting}[]
\FunctionTok{wget} \AttributeTok{{-}qO{-}}\NormalTok{ https://cloud.r{-}project.org/bin/linux/ubuntu/marutter\_pubkey.asc }\KeywordTok{|} \DataTypeTok{\textbackslash{}}
  \FunctionTok{sudo}\NormalTok{ tee }\AttributeTok{{-}a}\NormalTok{ /etc/apt/trusted.gpg.d/cran\_ubuntu\_key.asc}
\end{Highlighting}
\end{Shaded}

La huella digital (fingerprint) oficial de esta clave es:

\begin{verbatim}
E298A3A825C0D65DFD57CBB651716619E084DAB9
\end{verbatim}

Puedes verificarla con:

\begin{Shaded}
\begin{Highlighting}[]
\ExtensionTok{gpg} \AttributeTok{{-}{-}show{-}keys}\NormalTok{ /etc/apt/trusted.gpg.d/cran\_ubuntu\_key.asc}
\end{Highlighting}
\end{Shaded}

\textbf{Paso 3 --- Añadir el repositorio CRAN}

El comando usa \texttt{lsb\_release\ -cs} para detectar automáticamente
el nombre en clave de tu versión de Ubuntu (p.~ej. \texttt{noble},
\texttt{jammy}):

\begin{Shaded}
\begin{Highlighting}[]
\FunctionTok{sudo}\NormalTok{ add{-}apt{-}repository }\StringTok{"deb https://cloud.r{-}project.org/bin/linux/ubuntu }\VariableTok{$(}\ExtensionTok{lsb\_release} \AttributeTok{{-}cs}\VariableTok{)}\StringTok{{-}cran40/"}
\end{Highlighting}
\end{Shaded}

\begin{quote}
\texttt{cloud.r-project.org} redirige automáticamente a un espejo
(mirror) de CRAN cercano a tu ubicación.
\end{quote}

\textbf{Paso 4 --- Instalar R}

\begin{Shaded}
\begin{Highlighting}[]
\FunctionTok{sudo}\NormalTok{ apt update}
\FunctionTok{sudo}\NormalTok{ apt install }\AttributeTok{{-}{-}no{-}install{-}recommends}\NormalTok{ r{-}base}
\end{Highlighting}
\end{Shaded}

\textbf{Paso 5 --- Instalar herramientas de compilación (recomendado)}

Si vas a instalar paquetes de CRAN que requieren compilación desde
código fuente (lo cual es habitual con \texttt{install.packages()}),
instala también \texttt{r-base-dev}:

\begin{Shaded}
\begin{Highlighting}[]
\FunctionTok{sudo}\NormalTok{ apt install }\AttributeTok{{-}{-}no{-}install{-}recommends}\NormalTok{ r{-}base{-}dev}
\end{Highlighting}
\end{Shaded}

\subsection{2.3. Nota sobre versiones de Ubuntu
soportadas}\label{nota-sobre-versiones-de-ubuntu-soportadas}

CRAN publica binarios de R 4.x para las versiones de escritorio de
Ubuntu actualmente soportadas (hasta su fecha oficial de fin de vida).
Solo la última versión LTS recibe soporte completo. Si tu versión de
Kubuntu ya superó su fin de vida (EOL), considera actualizar el sistema
antes de continuar.

\subsection{\texorpdfstring{2.4. Instalación masiva de paquetes CRAN
como binarios
(\texttt{r2u})}{2.4. Instalación masiva de paquetes CRAN como binarios (r2u)}}\label{instalaciuxf3n-masiva-de-paquetes-cran-como-binarios-r2u}

Para evitar compilar cada paquete de R desde código fuente (proceso
lento), el proyecto \textbf{r2u}, mantenido por Dirk Eddelbuettel,
permite instalar miles de paquetes CRAN y Bioconductor como binarios
\texttt{.deb} listos para usar mediante un repositorio adicional. Es
opcional, pero muy recomendable si trabajas con muchos paquetes
(tidyverse, sf, etc.). Consulta la documentación del proyecto r2u para
los pasos de configuración específicos de tu versión de Ubuntu.

\begin{center}\rule{0.5\linewidth}{0.5pt}\end{center}

\section{3. Instalación de R en Arch
Linux}\label{instalaciuxf3n-de-r-en-arch-linux}

En Arch Linux, R está disponible directamente en el repositorio oficial
\textbf{extra}, gestionado por \texttt{pacman}. No se requieren
repositorios de terceros para obtener una versión reciente.

\subsection{3.1. Instalación con
pacman}\label{instalaciuxf3n-con-pacman}

\begin{Shaded}
\begin{Highlighting}[]
\FunctionTok{sudo}\NormalTok{ pacman }\AttributeTok{{-}Syu}\NormalTok{ r}
\end{Highlighting}
\end{Shaded}

Esto instala el paquete \texttt{r}, que incluye el binario \texttt{R},
\texttt{Rscript}, y las herramientas básicas de compilación de paquetes
de R (R en Arch ya integra lo equivalente a \texttt{r-base-dev}).

\subsection{3.2. Estructura de bibliotecas de paquetes en
Arch}\label{estructura-de-bibliotecas-de-paquetes-en-arch}

Según la documentación oficial (ArchWiki), conviene separar:

\begin{itemize}
\tightlist
\item
  \textbf{Paquetes gestionados por pacman}: se instalan en
  \texttt{/usr/lib/R/library}.
\item
  \textbf{Paquetes instalados por el usuario vía
  \texttt{install.packages()}}: se instalan en una ruta de usuario, por
  ejemplo \texttt{\textasciitilde{}/R/library} o similar, definida por
  \texttt{.libPaths()}.
\end{itemize}

La recomendación oficial es \textbf{no mezclar} ambos métodos de gestión
para un mismo paquete: si pacman gestiona un paquete de sistema, no debe
reinstalarse manualmente desde R, y viceversa.

\subsection{3.3. Paquetes de R adicionales vía
pacman}\label{paquetes-de-r-adicionales-vuxeda-pacman}

Arch empaqueta también una colección creciente de paquetes de R
individuales con el prefijo \texttt{r-} (por ejemplo
\texttt{r-tidyverse}, \texttt{r-ggplot2}, si están disponibles en los
repositorios oficiales o en AUR). Puedes buscarlos así:

\begin{Shaded}
\begin{Highlighting}[]
\ExtensionTok{pacman} \AttributeTok{{-}Ss}\NormalTok{ \^{}r{-}}
\end{Highlighting}
\end{Shaded}

\subsection{3.4. Instalación de paquetes CRAN desde dentro de
R}\label{instalaciuxf3n-de-paquetes-cran-desde-dentro-de-r}

Para paquetes no empaquetados por Arch, se instalan igual que en
cualquier sistema, desde la consola de R:

\begin{Shaded}
\begin{Highlighting}[]
\FunctionTok{install.packages}\NormalTok{(}\StringTok{"nombre\_del\_paquete"}\NormalTok{)}
\end{Highlighting}
\end{Shaded}

Si no tienes permisos de escritura en la biblioteca de sistema, R te
ofrecerá crear automáticamente una biblioteca personal en tu directorio
\texttt{\$HOME} (acepta con \texttt{yes}/\texttt{y} cuando se te
pregunte).

\subsection{3.5. Repositorios de terceros para binarios precompilados
(opcional)}\label{repositorios-de-terceros-para-binarios-precompilados-opcional}

Existen iniciativas comunitarias como \textbf{BioArchLinux} que ofrecen
binarios precompilados de paquetes CRAN/Bioconductor para Arch,
integrables con pacman mediante la herramienta \texttt{bspm} (permite
que \texttt{install.packages()} use pacman por debajo). Esto es opcional
y pensado para quienes instalan muchos paquetes con frecuencia; revisa
la wiki de BioArchLinux para la configuración exacta si te interesa.

\begin{center}\rule{0.5\linewidth}{0.5pt}\end{center}

\section{4. Instalación de R en
Windows}\label{instalaciuxf3n-de-r-en-windows}

\subsection{4.1. Descarga desde CRAN}\label{descarga-desde-cran}

\begin{enumerate}
\def\labelenumi{\arabic{enumi}.}
\tightlist
\item
  Visita el espejo oficial de CRAN: \textbf{https://cran.r-project.org/}
\item
  Haz clic en \textbf{``Download R for Windows''}.
\item
  Selecciona \textbf{``base''} (es la opción para una instalación
  nueva).
\item
  Haz clic en el enlace \textbf{``Download R-X.X.X for Windows''} (donde
  X.X.X es la versión más reciente).
\item
  Ejecuta el instalador \texttt{.exe} descargado.
\end{enumerate}

\subsection{4.2. Pasos del instalador}\label{pasos-del-instalador}

\begin{enumerate}
\def\labelenumi{\arabic{enumi}.}
\tightlist
\item
  Selecciona el idioma de instalación.
\item
  Acepta la licencia (GPL).
\item
  Elige la ruta de instalación (por defecto suele ser
  \texttt{C:\textbackslash{}Program\ Files\textbackslash{}R\textbackslash{}R-X.X.X}).
\item
  Selecciona los componentes a instalar (se recomienda dejar todos
  marcados: archivos principales de 64-bit, archivos de ayuda en HTML,
  traducciones de mensajes).
\item
  Configura las opciones de inicio (puedes dejar las opciones SDI/MDI
  por defecto si no tienes preferencia).
\item
  Elige si deseas crear un ícono en el escritorio y entradas en el menú
  de inicio.
\item
  Finaliza la instalación.
\end{enumerate}

\subsection{4.3. Verificación}\label{verificaciuxf3n}

Abre el menú de inicio y busca ``R x64 X.X.X'' o abre una terminal
(\texttt{cmd} o PowerShell) y ejecuta:

\begin{Shaded}
\begin{Highlighting}[]
\NormalTok{R {-}}\AttributeTok{{-}version}
\end{Highlighting}
\end{Shaded}

Si el comando no es reconocido, es porque R no se agregó automáticamente
al \texttt{PATH} del sistema; en ese caso, abre R desde su acceso
directo del menú de inicio, o agrega manualmente la carpeta \texttt{bin}
de la instalación (p.~ej.
\texttt{C:\textbackslash{}Program\ Files\textbackslash{}R\textbackslash{}R-X.X.X\textbackslash{}bin})
a la variable de entorno \texttt{PATH}.

\subsection{4.4. Rtools (opcional pero
recomendado)}\label{rtools-opcional-pero-recomendado}

Si vas a compilar paquetes de R desde código fuente en Windows
(necesario para algunos paquetes de CRAN que no tienen binario
precompilado para Windows), instala \textbf{Rtools}, disponible también
desde la página de descargas de CRAN para Windows (sección ``Rtools'').
Rtools provee las herramientas de compilación (gcc, make, etc.)
necesarias en ese entorno.

\begin{center}\rule{0.5\linewidth}{0.5pt}\end{center}

\section{5. Instalación de RStudio Desktop en Kubuntu /
Ubuntu}\label{instalaciuxf3n-de-rstudio-desktop-en-kubuntu-ubuntu}

\begin{quote}
\textbf{Prerrequisito}: R debe estar instalado antes de este paso (ver
sección 2).
\end{quote}

\subsection{\texorpdfstring{5.1. Descarga del paquete
\texttt{.deb}}{5.1. Descarga del paquete .deb}}\label{descarga-del-paquete-.deb}

\begin{enumerate}
\def\labelenumi{\arabic{enumi}.}
\tightlist
\item
  Visita la página oficial de descargas:
  \textbf{https://posit.co/download/rstudio-desktop/}
\item
  La página detecta tu sistema operativo y te ofrece el instalador
  correspondiente; descarga el paquete \texttt{.deb} para Ubuntu/Debian
  (arquitectura \texttt{amd64}).
\item
  Alternativamente, puedes descargar directamente con \texttt{wget}
  usando la URL del paquete vigente que indique la página de descargas
  (las URLs cambian con cada versión, por lo que se recomienda copiar el
  enlace actualizado directamente desde el sitio de Posit en el momento
  de instalar).
\end{enumerate}

\subsection{\texorpdfstring{5.2. Instalación del paquete
\texttt{.deb}}{5.2. Instalación del paquete .deb}}\label{instalaciuxf3n-del-paquete-.deb}

Desde la carpeta donde se descargó el archivo:

\begin{Shaded}
\begin{Highlighting}[]
\FunctionTok{sudo}\NormalTok{ apt install ./rstudio{-}2026.XX.X{-}XXX{-}amd64.deb}
\end{Highlighting}
\end{Shaded}

\begin{quote}
El uso de \texttt{apt\ install\ ./archivo.deb} (en vez de
\texttt{dpkg\ -i}) es preferible porque \texttt{apt} resuelve
automáticamente las dependencias faltantes.
\end{quote}

Si por alguna razón usaste \texttt{dpkg}, corrige dependencias faltantes
después con:

\begin{Shaded}
\begin{Highlighting}[]
\FunctionTok{sudo}\NormalTok{ dpkg }\AttributeTok{{-}i}\NormalTok{ rstudio{-}2026.XX.X{-}XXX{-}amd64.deb}
\FunctionTok{sudo}\NormalTok{ apt }\AttributeTok{{-}{-}fix{-}broken}\NormalTok{ install}
\end{Highlighting}
\end{Shaded}

\subsection{5.3. Verificación de la firma del paquete (opcional,
recomendado)}\label{verificaciuxf3n-de-la-firma-del-paquete-opcional-recomendado}

RStudio Desktop publica binarios firmados. Los usuarios de escritorio
Linux pueden necesitar importar la clave pública de firma de código de
Posit antes de instalar, dependiendo de la política de seguridad de tu
sistema. Esto es opcional para uso normal vía \texttt{apt\ install},
pero recomendable si quieres verificar la integridad del paquete
manualmente (especialmente si lo descargaste desde un espejo o enlace no
oficial).

\subsection{5.4. Distribuciones
soportadas}\label{distribuciones-soportadas}

RStudio Desktop genera binarios dirigidos a un conjunto específico de
distribuciones Linux de 64 bits. Revisa siempre la página oficial de
descargas para confirmar que tu versión de Kubuntu/Ubuntu está dentro de
las soportadas en el momento de tu instalación, ya que Posit retira
soporte para versiones antiguas con cada ciclo de lanzamientos (ver nota
de la sección 1.3 sobre Ubuntu 20.04).

\begin{center}\rule{0.5\linewidth}{0.5pt}\end{center}

\section{6. Instalación de RStudio Desktop en Arch
Linux}\label{instalaciuxf3n-de-rstudio-desktop-en-arch-linux}

\begin{quote}
\textbf{Importante}: RStudio Desktop \textbf{no está en los repositorios
oficiales} de Arch Linux (\texttt{core}/\texttt{extra}). Solo está
disponible mediante el \textbf{AUR} (Arch User Repository), que contiene
contenido producido por la comunidad, no por los desarrolladores
oficiales de Arch ni de Posit.
\end{quote}

\begin{quote}
\textbf{Prerrequisito}: R debe estar instalado antes de este paso (ver
sección 3).
\end{quote}

\subsection{6.1. Opciones disponibles en
AUR}\label{opciones-disponibles-en-aur}

Existen varios paquetes en AUR relacionados con RStudio; los dos
relevantes para RStudio Desktop (IDE de Posit) son:

\begin{ThreePartTable}

\begin{longtable}[]{@{}
  >{\raggedright\arraybackslash}p{(\linewidth - 6\tabcolsep) * \real{0.0833}}
  >{\raggedright\arraybackslash}p{(\linewidth - 6\tabcolsep) * \real{0.2609}}
  >{\raggedright\arraybackslash}p{(\linewidth - 6\tabcolsep) * \real{0.1630}}
  >{\raggedright\arraybackslash}p{(\linewidth - 6\tabcolsep) * \real{0.4928}}@{}}
\caption{Opciones disponibles en AUR para
RStudio}\label{tbl-aur-options}\tabularnewline
\toprule\noalign{}
\begin{minipage}[b]{\linewidth}\raggedright
Paquete AUR
\end{minipage} & \begin{minipage}[b]{\linewidth}\raggedright
Descripción
\end{minipage} & \begin{minipage}[b]{\linewidth}\raggedright
Ventaja
\end{minipage} & \begin{minipage}[b]{\linewidth}\raggedright
Desventaja
\end{minipage} \\
\midrule\noalign{}
\endfirsthead
\toprule\noalign{}
\begin{minipage}[b]{\linewidth}\raggedright
Paquete AUR
\end{minipage} & \begin{minipage}[b]{\linewidth}\raggedright
Descripción
\end{minipage} & \begin{minipage}[b]{\linewidth}\raggedright
Ventaja
\end{minipage} & \begin{minipage}[b]{\linewidth}\raggedright
Desventaja
\end{minipage} \\
\midrule\noalign{}
\endhead
\bottomrule\noalign{}
\endlastfoot
\texttt{rstudio-desktop-bin} & Empaqueta el \texttt{.deb} oficial de
Posit (binario distribuido para Ubuntu) & Instalación rápida, no
requiere compilar & Depende de que el mantenedor actualice el paquete
con cada release \\
\texttt{rstudio-desktop} & Compila RStudio desde el código fuente & Más
actualizado, optimizado para el sistema & Compilación larga (requiere
GWT, Boost, etc.) y más propensa a romperse con actualizaciones del
sistema (p.~ej. cambios en \texttt{libboost}) \\
\end{longtable}

\noindent \emph{Nota.} Elaboración~propia basada en la documentación de
AUR.

\end{ThreePartTable}

Para uso cotidiano sin compilar, \textbf{\texttt{rstudio-desktop-bin}}
es la opción más práctica para la mayoría de usuarios.

\subsection{6.2. Instalación con un ayudante de AUR
(recomendado)}\label{instalaciuxf3n-con-un-ayudante-de-aur-recomendado}

Si usas un ayudante de AUR como \texttt{yay} o \texttt{paru} (no
incluidos en Arch base, hay que instalarlos antes desde el propio AUR):

\begin{Shaded}
\begin{Highlighting}[]
\ExtensionTok{yay} \AttributeTok{{-}S}\NormalTok{ rstudio{-}desktop{-}bin}
\end{Highlighting}
\end{Shaded}

o, si prefieres compilar desde código fuente:

\begin{Shaded}
\begin{Highlighting}[]
\ExtensionTok{yay} \AttributeTok{{-}S}\NormalTok{ rstudio{-}desktop}
\end{Highlighting}
\end{Shaded}

\subsection{6.3. Instalación manual desde AUR (sin
ayudante)}\label{instalaciuxf3n-manual-desde-aur-sin-ayudante}

Si no usas un ayudante de AUR, el proceso oficial documentado por
ArchWiki es:

\begin{Shaded}
\begin{Highlighting}[]
\FunctionTok{git}\NormalTok{ clone https://aur.archlinux.org/rstudio{-}desktop{-}bin.git}
\BuiltInTok{cd}\NormalTok{ rstudio{-}desktop{-}bin}
\ExtensionTok{makepkg} \AttributeTok{{-}si}
\end{Highlighting}
\end{Shaded}

El flag \texttt{-si} le indica a \texttt{makepkg} que instale las
dependencias necesarias y luego instale el paquete resultante con
\texttt{pacman\ -U} automáticamente.

\subsection{6.4. Notas importantes para Arch /
Archcraft}\label{notas-importantes-para-arch-archcraft}

\begin{itemize}
\tightlist
\item
  El paquete \texttt{rstudio-desktop} (compilado desde fuente) puede
  romperse temporalmente tras actualizaciones de librerías base del
  sistema (por ejemplo, \texttt{libboost}); si esto ocurre, conviene
  forzar una reconstrucción con \texttt{pacman\ -Syudd} antes de
  reconstruir el paquete, o cambiar temporalmente a
  \texttt{rstudio-desktop-bin} mientras se resuelve.
\item
  Dado que ambos paquetes son mantenidos por voluntarios de la comunidad
  (no por Posit oficialmente), revisa los comentarios recientes en la
  página de AUR del paquete antes de instalar, por si hay incidencias
  activas no resueltas.
\item
  En sistemas con gestor de ventanas tipo \emph{tiling} tal como los
  usados en Archcraft, no debería haber incompatibilidades particulares
  con RStudio (es una aplicación Electron/Qt estándar), pero usuarios de
  Wayland han reportado necesitar banderas adicionales de lanzamiento
  para Electron; si tienes problemas de renderizado bajo Wayland, prueba
  forzando XWayland o revisa los comentarios del paquete en AUR para
  \emph{workarounds} conocidos.
\end{itemize}

\begin{center}\rule{0.5\linewidth}{0.5pt}\end{center}

\section{7. Instalación de RStudio Desktop en
Windows}\label{instalaciuxf3n-de-rstudio-desktop-en-windows}

\begin{quote}
\textbf{Prerrequisito}: R debe estar instalado antes de este paso (ver
sección 4).
\end{quote}

\subsection{7.1. Descarga}\label{descarga}

\begin{enumerate}
\def\labelenumi{\arabic{enumi}.}
\tightlist
\item
  Visita \textbf{https://posit.co/download/rstudio-desktop/}
\item
  Descarga el instalador autoextraíble (\texttt{.exe}) para Windows.
  RStudio para Windows requiere una edición de 64 bits de Windows 10 o
  superior.
\end{enumerate}

\subsection{7.2. Instalación}\label{instalaciuxf3n}

\begin{enumerate}
\def\labelenumi{\arabic{enumi}.}
\tightlist
\item
  Ejecuta el archivo \texttt{.exe} descargado.
\item
  Sigue el asistente de instalación (setup wizard).
\item
  Se requieren permisos de administrador para completar la instalación.
\item
  El instalador crea automáticamente un acceso directo de RStudio y una
  entrada en ``Agregar o quitar programas'' para una futura
  desinstalación.
\end{enumerate}

\subsection{7.3. Alternativa portable
(.zip)}\label{alternativa-portable-.zip}

Si prefieres no usar el instalador (por ejemplo, en un entorno
restringido o para tener una versión portable), Posit también distribuye
un archivo \texttt{.zip}:

\begin{enumerate}
\def\labelenumi{\arabic{enumi}.}
\tightlist
\item
  Descarga el \texttt{.zip} desde la misma página de descargas.
\item
  Extráelo en la ubicación que prefieras.
\item
  Ejecuta RStudio directamente desde \texttt{rstudio.exe}, ubicado en la
  subcarpeta \texttt{bin} del directorio extraído.
\end{enumerate}

\begin{center}\rule{0.5\linewidth}{0.5pt}\end{center}

\section{8. Verificación
post-instalación}\label{verificaciuxf3n-post-instalaciuxf3n}

\subsection{8.1. Verificar R desde terminal (Linux y
Windows)}\label{verificar-r-desde-terminal-linux-y-windows}

\begin{Shaded}
\begin{Highlighting}[]
\ExtensionTok{R} \AttributeTok{{-}{-}version}
\ExtensionTok{Rscript} \AttributeTok{{-}{-}version}
\end{Highlighting}
\end{Shaded}

Deberías ver información de versión similar a:

\begin{verbatim}
R version 4.X.X (20XX-XX-XX) -- "Nombre en clave"
\end{verbatim}

\subsection{8.2. Verificar R desde su propia consola
interactiva}\label{verificar-r-desde-su-propia-consola-interactiva}

\begin{Shaded}
\begin{Highlighting}[]
\ExtensionTok{R}
\end{Highlighting}
\end{Shaded}

Esto abre la consola interactiva de R. Para salir:

\begin{Shaded}
\begin{Highlighting}[]
\FunctionTok{q}\NormalTok{()}
\end{Highlighting}
\end{Shaded}

Te preguntará si deseas guardar la imagen del espacio de trabajo; puedes
responder \texttt{n} si no lo necesitas.

\subsection{8.3. Verificar que RStudio detecta
R}\label{verificar-que-rstudio-detecta-r}

Al abrir RStudio Desktop por primera vez, revisa la consola integrada
(panel inferior izquierdo por defecto): debe mostrar la versión de R
detectada al iniciar. Si tienes varias versiones de R instaladas, puedes
elegir cuál usar desde:

\texttt{Tools} → \texttt{Global\ Options} → \texttt{General} → sección
\textbf{R version} (en Linux/Windows el menú es equivalente, aunque la
ruta exacta puede variar ligeramente entre versiones de RStudio).

\begin{center}\rule{0.5\linewidth}{0.5pt}\end{center}

\section{9. Dependencias de sistema para paquetes de R
(Linux)}\label{dependencias-de-sistema-para-paquetes-de-r-linux}

Muchos paquetes de R populares (por ejemplo, los del \emph{tidyverse},
paquetes espaciales como \texttt{sf}, paquetes con gráficos avanzados,
etc.) requieren bibliotecas de sistema adicionales para compilarse
correctamente desde código fuente. Posit documenta un conjunto curado de
dependencias por distribución.

\subsection{9.1. Ejemplo para Ubuntu /
Debian}\label{ejemplo-para-ubuntu-debian}

A modo de referencia (lista no exhaustiva; revisa siempre la
documentación oficial de Posit para la lista completa y actualizada
según tu versión exacta de Ubuntu):

\begin{Shaded}
\begin{Highlighting}[]
\FunctionTok{sudo}\NormalTok{ apt update}
\FunctionTok{sudo}\NormalTok{ apt install }\AttributeTok{{-}y} \DataTypeTok{\textbackslash{}}
\NormalTok{  build{-}essential }\DataTypeTok{\textbackslash{}}
\NormalTok{  libcurl4{-}openssl{-}dev }\DataTypeTok{\textbackslash{}}
\NormalTok{  libssl{-}dev }\DataTypeTok{\textbackslash{}}
\NormalTok{  libxml2{-}dev }\DataTypeTok{\textbackslash{}}
\NormalTok{  libfontconfig1{-}dev }\DataTypeTok{\textbackslash{}}
\NormalTok{  libharfbuzz{-}dev }\DataTypeTok{\textbackslash{}}
\NormalTok{  libfribidi{-}dev }\DataTypeTok{\textbackslash{}}
\NormalTok{  libfreetype6{-}dev }\DataTypeTok{\textbackslash{}}
\NormalTok{  libpng{-}dev }\DataTypeTok{\textbackslash{}}
\NormalTok{  libtiff5{-}dev }\DataTypeTok{\textbackslash{}}
\NormalTok{  libjpeg{-}dev }\DataTypeTok{\textbackslash{}}
\NormalTok{  libgit2{-}dev}
\end{Highlighting}
\end{Shaded}

Esta lista cubre dependencias habituales de paquetes como
\texttt{devtools}, \texttt{httr}, \texttt{xml2}, \texttt{ragg},
\texttt{textshaping}, \texttt{git2r}, entre otros.

\subsection{9.2. Ejemplo para Arch Linux}\label{ejemplo-para-arch-linux}

En Arch, la mayoría de estas bibliotecas suelen instalarse
automáticamente como dependencias declaradas en el \texttt{PKGBUILD} del
paquete \texttt{r-*} correspondiente o se resuelven directamente al
compilar con \texttt{install.packages()}, siempre que tengas el grupo
\texttt{base-devel} instalado:

\begin{Shaded}
\begin{Highlighting}[]
\FunctionTok{sudo}\NormalTok{ pacman }\AttributeTok{{-}S} \AttributeTok{{-}{-}needed}\NormalTok{ base{-}devel}
\end{Highlighting}
\end{Shaded}

\subsection{9.3. Detección automática de
dependencias}\label{detecciuxf3n-automuxe1tica-de-dependencias}

Posit ofrece una herramienta de \textbf{detección de dependencias de
sistema} que analiza qué bibliotecas de sistema requiere un conjunto de
paquetes de R antes de instalarlos, evitando errores de compilación. Es
útil si planeas instalar muchos paquetes de una sola vez (por ejemplo,
en un servidor o máquina nueva). Consulta la documentación oficial de
Posit (``System Dependency Detection'') para más detalles sobre su uso.

\begin{center}\rule{0.5\linewidth}{0.5pt}\end{center}

\section{10. Gestión de paquetes de
R}\label{gestiuxf3n-de-paquetes-de-r}

\subsection{10.1. Instalar paquetes desde
CRAN}\label{instalar-paquetes-desde-cran}

Desde la consola de R (en cualquier sistema operativo):

\begin{Shaded}
\begin{Highlighting}[]
\FunctionTok{install.packages}\NormalTok{(}\StringTok{"ggplot2"}\NormalTok{)}
\FunctionTok{install.packages}\NormalTok{(}\FunctionTok{c}\NormalTok{(}\StringTok{"dplyr"}\NormalTok{, }\StringTok{"tidyr"}\NormalTok{, }\StringTok{"readr"}\NormalTok{))  }\CommentTok{\# instalación múltiple}
\end{Highlighting}
\end{Shaded}

\subsection{10.2. Cargar un paquete
instalado}\label{cargar-un-paquete-instalado}

\begin{Shaded}
\begin{Highlighting}[]
\FunctionTok{library}\NormalTok{(ggplot2)}
\end{Highlighting}
\end{Shaded}

\subsection{10.3. Actualizar todos los paquetes
instalados}\label{actualizar-todos-los-paquetes-instalados}

\begin{Shaded}
\begin{Highlighting}[]
\FunctionTok{update.packages}\NormalTok{(}\AttributeTok{ask =} \ConstantTok{FALSE}\NormalTok{)}
\end{Highlighting}
\end{Shaded}

\subsection{10.4. Ver la biblioteca de paquetes
activa}\label{ver-la-biblioteca-de-paquetes-activa}

\begin{Shaded}
\begin{Highlighting}[]
\FunctionTok{.libPaths}\NormalTok{()}
\end{Highlighting}
\end{Shaded}

\subsection{\texorpdfstring{10.5. Instalar paquetes desde GitHub
(requiere \texttt{remotes} o
\texttt{devtools})}{10.5. Instalar paquetes desde GitHub (requiere remotes o devtools)}}\label{instalar-paquetes-desde-github-requiere-remotes-o-devtools}

\begin{Shaded}
\begin{Highlighting}[]
\FunctionTok{install.packages}\NormalTok{(}\StringTok{"remotes"}\NormalTok{)}
\NormalTok{remotes}\SpecialCharTok{::}\FunctionTok{install\_github}\NormalTok{(}\StringTok{"usuario/repositorio"}\NormalTok{)}
\end{Highlighting}
\end{Shaded}

\begin{center}\rule{0.5\linewidth}{0.5pt}\end{center}

\section{11. Actualización y
desinstalación}\label{actualizaciuxf3n-y-desinstalaciuxf3n}

\subsection{11.1. Actualizar R en Kubuntu/Ubuntu (repositorio
CRAN)}\label{actualizar-r-en-kubuntuubuntu-repositorio-cran}

\begin{Shaded}
\begin{Highlighting}[]
\FunctionTok{sudo}\NormalTok{ apt update}
\FunctionTok{sudo}\NormalTok{ apt upgrade r{-}base}
\end{Highlighting}
\end{Shaded}

\subsection{11.2. Actualizar R en Arch
Linux}\label{actualizar-r-en-arch-linux}

Arch usa un modelo de lanzamiento continuo (\emph{rolling release}); R
se actualiza junto con el resto del sistema:

\begin{Shaded}
\begin{Highlighting}[]
\FunctionTok{sudo}\NormalTok{ pacman }\AttributeTok{{-}Syu}
\end{Highlighting}
\end{Shaded}

\subsection{11.3. Actualizar RStudio
Desktop}\label{actualizar-rstudio-desktop}

En todos los sistemas, la actualización de RStudio Desktop se realiza
descargando e instalando el nuevo paquete/instalador desde la página
oficial de descargas (no existe, en la edición open-source, un mecanismo
de auto-actualización integrado en la aplicación). En Arch vía AUR,
basta con actualizar el paquete con tu ayudante de AUR habitual
(\texttt{yay\ -Syu}, por ejemplo).

\subsection{11.4. Desinstalación}\label{desinstalaciuxf3n}

\textbf{Kubuntu/Ubuntu:}

\begin{Shaded}
\begin{Highlighting}[]
\FunctionTok{sudo}\NormalTok{ apt remove r{-}base r{-}base{-}dev}
\FunctionTok{sudo}\NormalTok{ apt remove rstudio}
\end{Highlighting}
\end{Shaded}

\textbf{Arch Linux:}

\begin{Shaded}
\begin{Highlighting}[]
\FunctionTok{sudo}\NormalTok{ pacman }\AttributeTok{{-}Rns}\NormalTok{ r}
\ExtensionTok{yay} \AttributeTok{{-}Rns}\NormalTok{ rstudio{-}desktop{-}bin   }\CommentTok{\# o rstudio{-}desktop, según cuál hayas instalado}
\end{Highlighting}
\end{Shaded}

\textbf{Windows:}

Usa el panel ``Agregar o quitar programas'' de Windows y desinstala
``R'' y ``RStudio'' por separado, ya que son aplicaciones
independientes.

\begin{center}\rule{0.5\linewidth}{0.5pt}\end{center}

\section{12. Solución de problemas
comunes}\label{soluciuxf3n-de-problemas-comunes}

\subsection{\texorpdfstring{12.1.
\texttt{add-apt-repository:\ command\ not\ found}
(Ubuntu/Kubuntu)}{12.1. add-apt-repository: command not found (Ubuntu/Kubuntu)}}\label{add-apt-repository-command-not-found-ubuntukubuntu}

Falta el paquete \texttt{software-properties-common}. Instálalo con:

\begin{Shaded}
\begin{Highlighting}[]
\FunctionTok{sudo}\NormalTok{ apt install software{-}properties{-}common}
\end{Highlighting}
\end{Shaded}

\subsection{12.2. Error de clave GPG no válida al añadir el repositorio
CRAN}\label{error-de-clave-gpg-no-vuxe1lida-al-auxf1adir-el-repositorio-cran}

Verifica que la huella digital coincide con
\texttt{E298A3A825C0D65DFD57CBB651716619E084DAB9}. Si la descarga de la
clave falla, intenta nuevamente o usa un espejo de CRAN alternativo
listado en https://cran.r-project.org/mirrors.html.

\subsection{12.3. RStudio no detecta R recién instalado
(Linux)}\label{rstudio-no-detecta-r-reciuxe9n-instalado-linux}

Cierra y vuelve a abrir RStudio. Si el problema persiste, revisa
\texttt{Tools} → \texttt{Global\ Options} → \texttt{General} y
selecciona manualmente la ruta del binario de R (normalmente
\texttt{/usr/lib/R} en sistemas basados en Debian/Ubuntu, o
\texttt{/usr/lib/R} también en Arch tras instalar el paquete
\texttt{r}).

\subsection{12.4. Falla la compilación de un paquete de R por falta de
una librería de sistema
(Linux)}\label{falla-la-compilaciuxf3n-de-un-paquete-de-r-por-falta-de-una-libreruxeda-de-sistema-linux}

El mensaje de error en R suele indicar el nombre de la librería faltante
(por ejemplo, \texttt{libcurl} o \texttt{libxml2}). Instala el paquete
de desarrollo (\texttt{-dev} en Debian/Ubuntu, sin sufijo o con nombre
equivalente en Arch) correspondiente y reintenta
\texttt{install.packages()}.

\subsection{\texorpdfstring{12.5. El paquete \texttt{rstudio-desktop} de
AUR no compila tras una actualización del sistema
(Arch)}{12.5. El paquete rstudio-desktop de AUR no compila tras una actualización del sistema (Arch)}}\label{el-paquete-rstudio-desktop-de-aur-no-compila-tras-una-actualizaciuxf3n-del-sistema-arch}

Es un problema conocido y documentado en los comentarios del propio
paquete AUR, generalmente vinculado a cambios en \texttt{libboost} u
otras librerías del sistema. Prueba:

\begin{Shaded}
\begin{Highlighting}[]
\FunctionTok{sudo}\NormalTok{ pacman }\AttributeTok{{-}Syudd}
\end{Highlighting}
\end{Shaded}

antes de reconstruir, o cambia temporalmente al paquete
\texttt{rstudio-desktop-bin}, que usa un binario precompilado y no
depende de compilar contra las librerías actuales del sistema.

\subsection{\texorpdfstring{12.6. En Windows, \texttt{R} no es
reconocido como
comando}{12.6. En Windows, R no es reconocido como comando}}\label{en-windows-r-no-es-reconocido-como-comando}

El instalador de R para Windows no siempre agrega el ejecutable al
\texttt{PATH} del sistema automáticamente. Agrega manualmente la ruta
\texttt{bin} de tu instalación de R (p.~ej.
\texttt{C:\textbackslash{}Program\ Files\textbackslash{}R\textbackslash{}R-X.X.X\textbackslash{}bin\textbackslash{}x64})
a la variable de entorno \texttt{PATH} desde el Panel de Control →
Sistema → Configuración avanzada del sistema → Variables de entorno.

\begin{center}\rule{0.5\linewidth}{0.5pt}\end{center}

\section{13. Referencias oficiales}\label{referencias-oficiales}

\begin{itemize}
\tightlist
\item
  \textbf{CRAN --- R Project}: https://cran.r-project.org/
\item
  \textbf{CRAN --- Instrucciones para Ubuntu (breve)}:
  https://cran.r-project.org/bin/linux/ubuntu/
\item
  \textbf{CRAN --- Instrucciones para Ubuntu (completas / README)}:
  https://cran.r-project.org/bin/linux/ubuntu/fullREADME.html
\item
  \textbf{CRAN --- R para Windows}:
  https://cran.r-project.org/bin/windows/base/
\item
  \textbf{ArchWiki --- R}: https://wiki.archlinux.org/title/R
\item
  \textbf{ArchWiki --- Pacman}: https://wiki.archlinux.org/title/Pacman
\item
  \textbf{AUR --- rstudio-desktop-bin}:
  https://aur.archlinux.org/packages/rstudio-desktop-bin
\item
  \textbf{AUR --- rstudio-desktop}:
  https://aur.archlinux.org/packages/rstudio-desktop
\item
  \textbf{Posit --- Descargas de RStudio Desktop}:
  https://posit.co/download/rstudio-desktop/
\item
  \textbf{Posit Docs --- Instalación de R (Install R)}:
  https://docs.posit.co/resources/install-r.html
\item
  \textbf{Posit Docs --- Soporte de plataformas}:
  https://docs.posit.co/platform-support.html
\item
  \textbf{Posit Docs --- Guía del IDE de RStudio}:
  https://docs.posit.co/ide/user/
\end{itemize}

\begin{center}\rule{0.5\linewidth}{0.5pt}\end{center}

\section{Publicaciones Similares}\label{publicaciones-similares}

Si te interesó este artículo, te recomendamos que explores otros blogs y
recursos relacionados que pueden ampliar tus conocimientos. Aquí te dejo
algunas sugerencias:

\begin{enumerate}
\def\labelenumi{\arabic{enumi}.}
\tightlist
\item
  \href{https://numerus-scriptum.netlify.app/r/2020-06-10-011-instalacion-de-r/index.pdf}{\faIcon{file-pdf}}
  \href{https://numerus-scriptum.netlify.app/r/2020-06-10-011-instalacion-de-r}{011
  Instalacion De R}
\item
  \href{https://numerus-scriptum.netlify.app/r/2020-06-10-012-configurar-entorno-virtual-r/index.pdf}{\faIcon{file-pdf}}
  \href{https://numerus-scriptum.netlify.app/r/2020-06-10-012-configurar-entorno-virtual-r}{012
  Configurar Entorno Virtual R}
\item
  \href{https://numerus-scriptum.netlify.app/r/2020-06-10-013-lo-que-debemos-saber-de-r/index.pdf}{\faIcon{file-pdf}}
  \href{https://numerus-scriptum.netlify.app/r/2020-06-10-013-lo-que-debemos-saber-de-r}{013
  Lo Que Debemos Saber De R}
\item
  \href{https://numerus-scriptum.netlify.app/r/2021-04-05-02-manipulacion-de-datos/index.pdf}{\faIcon{file-pdf}}
  \href{https://numerus-scriptum.netlify.app/r/2021-04-05-02-manipulacion-de-datos}{02
  Manipulacion De Datos}
\item
  \href{https://numerus-scriptum.netlify.app/r/2021-04-12-03-visualizacion-de-datos/index.pdf}{\faIcon{file-pdf}}
  \href{https://numerus-scriptum.netlify.app/r/2021-04-12-03-visualizacion-de-datos}{03
  Visualizacion De Datos}
\item
  \href{https://numerus-scriptum.netlify.app/r/2022-11-07-04-modelo-de-machine-learning-i-analisis-exploratorio/index.pdf}{\faIcon{file-pdf}}
  \href{https://numerus-scriptum.netlify.app/r/2022-11-07-04-modelo-de-machine-learning-i-analisis-exploratorio}{04
  Modelo De Machine Learning I Analisis Exploratorio}
\item
  \href{https://numerus-scriptum.netlify.app/r/2022-11-14-05-modelo-de-machine-learning-ii-modelo-de-clasificacion/index.pdf}{\faIcon{file-pdf}}
  \href{https://numerus-scriptum.netlify.app/r/2022-11-14-05-modelo-de-machine-learning-ii-modelo-de-clasificacion}{05
  Modelo De Machine Learning Ii Modelo De Clasificacion}
\item
  \href{https://numerus-scriptum.netlify.app/r/2022-11-21-06-modelo-de-machine-learning-iii-modelo-de-regresion/index.pdf}{\faIcon{file-pdf}}
  \href{https://numerus-scriptum.netlify.app/r/2022-11-21-06-modelo-de-machine-learning-iii-modelo-de-regresion}{06
  Modelo De Machine Learning Iii Modelo De Regresion}
\item
  \href{https://numerus-scriptum.netlify.app/r/2022-11-28-07-modelo-de-machine-learning-iv-tex-mining/index.pdf}{\faIcon{file-pdf}}
  \href{https://numerus-scriptum.netlify.app/r/2022-11-28-07-modelo-de-machine-learning-iv-tex-mining}{07
  Modelo De Machine Learning Iv Tex Mining}
\end{enumerate}

Esperamos que encuentres estas publicaciones igualmente interesantes y
útiles. ¡Disfruta de la lectura!






\end{document}
